% AdvancedEMQuality20261023
\section{静磁场}

\subsection{电流分布、矢势与磁能}

稳恒电流满足 $\nabla\cdot\vect j=0$，静磁场方程为
\[
 \nabla\times\vect B=\mu_0\vect j,
 \qquad \nabla\cdot\vect B=0.
\]
因此可引入矢势
\[
 \boxed{\vect A(\vect r)=\frac{\mu_0}{4\pi}
 \int\frac{\vect j(\vect r')}{|\vect r-\vect r'|}\dd V'},
 \qquad \vect B=\nabla\times\vect A,
\]
即 Biot--Savart 公式
\[
 \vect B=\frac{\mu_0}{4\pi}\int
 \frac{\vect j(\vect r')\times(\vect r-\vect r')}
 {|\vect r-\vect r'|^3}\dd V'.
\]
对线电流，$\vect j\dd V$ 替换为 $I\dd\vect l$。积分形式为
\[
 \oint\vect B\cdot\dd\vect S=0,
 \qquad \oint\vect B\cdot\dd\vect l=\mu_0I_{\rm enc},
 \qquad \int_S\vect B\cdot\dd\vect S=\oint_{\partial S}\vect A\cdot\dd\vect l.
\]

\begin{advexamplebox}{Biot--Savart 矢势的 Coulomb 规范}\label{advA4:ex:coulomb-gauge}
证明在稳恒局域电流条件下，上述 $\vect A$ 自动满足
$\nabla\cdot\vect A=0$。

\begin{center}\begin{minipage}{.96\linewidth}\centering
\begin{tikzpicture}[>=Latex,font=\small,scale=.85,every node/.style={fill=none}]

\draw[thick] (0,0) ellipse (1.2 and .6);\draw[DeepBlue,->] (1.2,0) arc[start angle=0,end angle=65,x radius=1.2,y radius=.6];
\fill (2.5,1.2) circle (1.7pt) node[right]{$\boldsymbol r$};\fill (0,.6) circle (1.5pt) node[above]{$\boldsymbol r'$};\draw[dashed] (0,.6)--(2.5,1.2) node[midway,above]{$|\boldsymbol r-\boldsymbol r'|$};
\node[below] at (0,-.9){局域稳恒电流，$\nabla'\cdot\boldsymbol j=0$};

\end{tikzpicture}
\figcap{题图：源点与场点的导数作用在同一个距离核上，源电流局域且无散。}\label{advA4:quality:A4-magnetostatics:q1}
\end{minipage}\end{center}
\end{advexamplebox}
\global\AdvHideHeadingtrue
\begin{advsolution}
\begin{center}\begin{minipage}{.96\linewidth}\centering
% EMAdvFigureHeading
\raggedright{\bfseries\color{DeepBlue}解答}\par\smallskip\centering

\begin{tikzpicture}[>=Latex,font=\small,scale=.85,every node/.style={fill=none}]

\draw[thick] (0,0) ellipse (1.2 and .6);\draw[DeepBlue,->] (1.2,0) arc[start angle=0,end angle=65,x radius=1.2,y radius=.6];
\fill (2.5,1.2) circle (1.7pt) node[right]{$\boldsymbol r$};\fill (0,.6) circle (1.5pt) node[above]{$\boldsymbol r'$};\draw[dashed] (0,.6)--(2.5,1.2) node[midway,above]{$|\boldsymbol r-\boldsymbol r'|$};
\node[below] at (0,-.9){局域稳恒电流，$\nabla'\cdot\boldsymbol j=0$};

\end{tikzpicture}
\figcap{解图：源点与场点的导数作用在同一个距离核上，源电流局域且无散。}\label{advA4:quality:A4-magnetostatics:a1}
\end{minipage}\end{center}

利用 $\nabla(1/R)=-\nabla'(1/R)$，
\[
 \nabla\cdot\vect A
 =\frac{\mu_0}{4\pi}\int\vect j'\cdot\nabla\frac1R\dd V'
 =-\frac{\mu_0}{4\pi}\int\vect j'\cdot\nabla'\frac1R\dd V'.
\]
把被积式写成
\[
 \vect j'\cdot\nabla'\frac1R
 =\nabla'\cdot\frac{\vect j'}R-\frac{\nabla'\cdot\vect j'}R.
\]
第二项因稳恒连续性 $\nabla'\cdot\vect j'=0$ 消失；第一项化为无穷远边界通量，局域电流使其为零。因此
\[
 \boxed{\nabla\cdot\vect A=0}.
\]
若电流穿过所选有限区域边界，边界项不能任意丢弃；应把回流路径一并计入。

图~\ref{advA4:quality:A4-magnetostatics:a1} 中，同一核满足 $\nabla(1/R)=-\nabla'(1/R)$；对源坐标分部积分时，体项由稳恒连续性为零，远处面项由局域性为零。
\end{advsolution}

静磁能为
\[
 U_B=\frac12\int\vect j\cdot\vect A\dd V.
\]

\begin{advtheorem}{静磁 Green 互易关系}\label{advA4:thm:magnetic-reciprocity}
两组局域稳恒电流及其 Coulomb 规范矢势满足
\[
 \boxed{\int\vect j_1\cdot\vect A_2\dd V
 =\int\vect j_2\cdot\vect A_1\dd V}.
\]
\end{advtheorem}
\noindent\textbf{证明}\quad 代入积分核，
\[
 \int\vect j_1(\vect r)\cdot\vect A_2(\vect r)\dd V
 =\frac{\mu_0}{4\pi}\iint
 \frac{\vect j_1(\vect r)\cdot\vect j_2(\vect r')}
 {|\vect r-\vect r'|}\dd V\dd V'.
\]
交换 $\vect r,\vect r'$ 及下标1、2，积分不变，即得结论。\hfill$\square$

\begin{advtheorem}{磁场能}\label{advA4:thm:magnetic-energy}
局域稳恒电流的能量也可写成
\[
 \boxed{\frac12\int\vect j\cdot\vect A\dd V
 =\frac1{2\mu_0}\int B^2\dd V}.
\]
\end{advtheorem}
\noindent\textbf{证明}\quad 由 $\vect j=\mu_0^{-1}\nabla\times\vect B$ 和
\[
 \nabla\cdot(\vect A\times\vect B)
 =\vect B\cdot\nabla\times\vect A-
 \vect A\cdot\nabla\times\vect B
\]
得到
\[
 \int\vect j\cdot\vect A\dd V
 =\frac1{\mu_0}\int B^2\dd V
 -\frac1{\mu_0}\oint(\vect A\times\vect B)\cdot\dd\vect S.
\]
局域场在无穷远的边界项消失，乘以 $1/2$ 即得。\hfill$\square$

\begin{advexamplebox}{线圈系统的电感矩阵}\label{advA4:ex:inductance-matrix}
证明真空或线性互易介质中 $N$ 个线圈的磁能为
\[
 U_B=\frac12\sum_{i,j}L_{ij}I_iI_j,
\]
其中穿过线圈 $i$ 的磁链 $\Phi_i=\sum_jL_{ij}I_j$，并证明
$L_{ij}=L_{ji}$。说明固定电流求机械力时为什么使用磁共能。

\begin{center}\begin{minipage}{.96\linewidth}\centering
\begin{tikzpicture}[>=Latex,font=\small,scale=.85,every node/.style={fill=none}]

\draw[thick] (0,0) ellipse (1.1 and .5);\draw[thick] (2.6,.7) ellipse (1.1 and .5);\node at (0,-.8){线圈 $i,I_i$};\node at (2.6,1.6){线圈 $j,I_j$};\draw[dashed] (.5,.35)--(2.1,.45) node[midway,above]{$M_{ij}$};

\end{tikzpicture}
\figcap{题图：相互交链的线圈系统，各电流和磁链都带同一个端口方向约定。}\label{advA4:quality:A4-magnetostatics:q2}
\end{minipage}\end{center}
\end{advexamplebox}
\global\AdvHideHeadingtrue
\begin{advsolution}
\begin{center}\begin{minipage}{.96\linewidth}\centering
% EMAdvFigureHeading
\raggedright{\bfseries\color{DeepBlue}解答}\par\smallskip\centering

\begin{tikzpicture}[>=Latex,font=\small,scale=.85,every node/.style={fill=none}]

\draw[thick] (0,0) ellipse (1.1 and .5);\draw[thick] (2.6,.7) ellipse (1.1 and .5);\node at (0,-.8){线圈 $i,I_i$};\node at (2.6,1.6){线圈 $j,I_j$};\draw[dashed] (.5,.35)--(2.1,.45) node[midway,above]{$M_{ij}$};

\end{tikzpicture}
\figcap{解图：相互交链的线圈系统，各电流和磁链都带同一个端口方向约定。}\label{advA4:quality:A4-magnetostatics:a2}
\end{minipage}\end{center}

% EMReciprocalMedium
“互易”是这里不可省略的条件。真空可直接使用前述积分核；对于位置可变但满足 $\mu_{ij}=\mu_{ji}$ 的线性磁介质，忽略无穷远面项后有
\[
 \int\boldsymbol j_1\cdot\boldsymbol A_2\,\mathrm dV
 =\int\boldsymbol H_1\cdot\boldsymbol B_2\,\mathrm dV
 =\int\boldsymbol H_2\cdot\boldsymbol B_1\,\mathrm dV
 =\int\boldsymbol j_2\cdot\boldsymbol A_1\,\mathrm dV.
\]
中间等号才使用介质张量的对称性。一般非互易介质不能仅凭“线性”就声称电感矩阵对称。

把总电流写成 $\vect j=\sum_iI_i\vect j_i^{(1)}$，矢势也线性叠加。于是
\[
 U_B=\frac12\sum_{i,j}I_iI_j
 \int\vect j_i^{(1)}\cdot\vect A_j^{(1)}\dd V
 \equiv\frac12\sum_{i,j}L_{ij}I_iI_j.
\]
对细线圈，$\int\vect j_i^{(1)}\cdot\vect A_j^{(1)}\dd V
=\oint_i\vect A_j^{(1)}\cdot\dd\vect l=\Phi_{i\leftarrow j}/I_j$，故
$\Phi_i=\partial U_B/\partial I_i=\sum_jL_{ij}I_j$。定理\ref{advA4:thm:magnetic-reciprocity}立即给出
$L_{ij}=L_{ji}$。

若外电源固定 $I_i$，构型改变时电源还会做功
$\sum_iI_i\dd\Phi_i$。Legendre 变换后的机械势为
\[
 V_{\rm mech}=U_B-\sum_iI_i\Phi_i=-U_B
\]
（线性体系），所以广义力
$Q_\alpha=\partial U_B/\partial q_\alpha|_I$。若超导回路固定磁链，则直接在固定 $\Phi$ 下用场能，两个条件不能混用。

图~\ref{advA4:quality:A4-magnetostatics:a2} 中，互感项连接两个不同线圈，能量双重求和中的 $1/2$ 防止把同一对耦合重复计数；电流固定的机械变化必须把电源做功一起核算。
\end{advsolution}

\subsection{磁多极矩}

远区矢势可按
$1/|\vect r-\vect r'|$ 展开。零阶项含
$\int\vect j\dd V=0$；其证明是对任意常矢量 $\vect c$ 写
\[
 \vect c\cdot\int\vect j\dd V
 =\int\nabla\cdot[(\vect c\cdot\vect r)\vect j]\dd V=0.
\]
一阶矩张量为
\[
 M_{ij}=\int x_i j_j\dd V.
\]

\begin{advtheorem}{稳恒电流的一阶矩反对称}\label{advA4:thm:moment-antisymmetric}
\[
 \boxed{M_{ij}+M_{ji}=0}.
\]
因而可用磁偶极矩
\[
 \boxed{\vect m=\frac12\int\vect r\times\vect j\dd V}
\]
完全表示它。
\end{advtheorem}
\noindent\textbf{证明}\quad 对局域电流积分
\[
 \partial_k(x_ix_jj_k)=x_jj_i+x_ij_j+x_ix_j\partial_kj_k.
\]
最后一项由 $\nabla\cdot\vect j=0$ 消失；左边体积分化为无穷远边界通量并为零，故
$\int(x_jj_i+x_ij_j)\dd V=0$。反对称二阶张量在三维与矢量对偶，正是所列 $\vect m$。\hfill$\square$

远区矢势和磁场为
\[
 \vect A=\frac{\mu_0}{4\pi}\frac{\vect m\times\vect r}{r^3},
 \qquad (r\ne0).
\]

\begin{advtheorem}{点磁偶极子的完整分布场}\label{advA4:thm:magnetic-dipole-field}
\[
 \boxed{\vect B(\vect r)=\frac{\mu_0}{4\pi}
 \frac{3(\vect m\cdot\vect r)\vect r-r^2\vect m}{r^5}
 +\frac{2\mu_0}{3}\vect m\,\delta^{(3)}(\vect r)}.
\]
\end{advtheorem}
\noindent\textbf{证明}\quad 由
$\vect A=-(\mu_0/4\pi)\vect m\times\nabla(1/r)$ 计算旋度，并使用分布 Hessian
\[
 \partial_i\partial_j\frac1r
 =\frac{3x_ix_j-r^2\delta_{ij}}{r^5}
 -\frac{4\pi}{3}\delta_{ij}\delta^{(3)}(\vect r).
\]
两个缩并项组合后，$\delta$ 系数为 $2\mu_0/3$。该项保证
$\nabla\cdot\vect B=0$ 在原点也成立。\hfill$\square$

外场缓变时
\[
 \vect M=\vect m\times\vect B,
 \qquad \vect F=\nabla(\vect m\cdot\vect B),
\]
可由小电流环或势能 $U=-\vect m\cdot\vect B$ 推得。

\begin{advexamplebox}{可伸缩超导螺线管}\label{advA4:ex:superconducting-solenoid}
轻质软螺线管截面积恒为 $S$、匝数 $N$、初长 $l$、初始电流 $I$。两端各接质量 $m$ 的质点，初态静止，释放后求长度 $x(t)$ 和电流 $i(t)$，忽略边缘效应。

\begin{center}
\begin{tikzpicture}[x=1cm,y=1cm,>=Latex]
 \draw[thick] (-3,0) circle(.32) node {$m$};
 \draw[thick] (3,0) circle(.32) node {$m$};
 \draw[thick] (-2.68,0) -- (-2.35,0);
 \draw[thick,decorate,decoration={coil,segment length=5.5pt,amplitude=5pt}]
   (-2.35,0)--(2.35,0);
 \draw[thick] (2.35,0)--(2.68,0);
 \draw[<->] (-2.35,-.72)--(2.35,-.72) node[midway,fill=none,below] {$x(t)$};
 \draw[->,DeepBlue] (-3.45,0)--(-3.9,0) node[left] {$+q_L$};
 \draw[->,DeepBlue] (3.45,0)--(3.9,0) node[right] {$+q_R$};
 \node[above] at (0,.48) {$N$ 匝，截面积 $S$};
\node at(0,-1.5){箭头规定两端向外位移的正向，非实际速度};
\end{tikzpicture}
\figcap{题图：端点质点、可变线圈长度和固定截面积，原图保留。}\label{advA4:quality:A4-magnetostatics:q3}

\end{center}
\end{advexamplebox}
\global\AdvHideHeadingtrue
\begin{advsolution}
\begin{center}\begin{minipage}{.96\linewidth}\centering
% EMAdvFigureHeading
\raggedright{\bfseries\color{DeepBlue}解答}\par\smallskip\centering

\begin{tikzpicture}[x=1cm,y=1cm,>=Latex]
 \draw[thick] (-3,0) circle(.32) node {$m$};
 \draw[thick] (3,0) circle(.32) node {$m$};
 \draw[thick] (-2.68,0) -- (-2.35,0);
 \draw[thick,decorate,decoration={coil,segment length=5.5pt,amplitude=5pt}]
   (-2.35,0)--(2.35,0);
 \draw[thick] (2.35,0)--(2.68,0);
 \draw[<->] (-2.35,-.72)--(2.35,-.72) node[midway,fill=none,below] {$x(t)$};
 \draw[->,DeepBlue] (-3.45,0)--(-3.9,0) node[left] {$+q_L$};
 \draw[->,DeepBlue] (3.45,0)--(3.9,0) node[right] {$+q_R$};
 \node[above] at (0,.48) {$N$ 匝，截面积 $S$};
\node at(0,-1.5){箭头规定两端向外位移的正向，非实际速度};
\end{tikzpicture}
\figcap{解图：端点质点、可变线圈长度和固定截面积，原图保留。}\label{advA4:quality:A4-magnetostatics:a3}
\end{minipage}\end{center}

长螺线管自感
\[
 L(x)=\frac{\mu_0N^2S}{x}.
\]
超导回路磁链守恒：
\[
 L(x)i(t)=L(l)I
 \quad\Longrightarrow\quad
 \boxed{i(t)=I\frac{x(t)}l}.
\]
固定磁链下的磁能
\[
 U_B=\frac12L(x)i^2
 =\frac{\mu_0N^2SI^2}{2l^2}x\equiv ax.
\]
两质点关于质心对称运动，各自速度大小 $|\dot x|/2$，总动能
$K=m\dot x^2/4$。能量守恒
\[
 \frac m4\dot x^2+ax=al
\]
或约化质量方程 $(m/2)\ddot x=-a$ 给出
\[
 \boxed{x(t)=l-\frac{\mu_0N^2SI^2}{2ml^2}t^2},
\]
\[
 \boxed{i(t)=I\left(1-\frac{\mu_0N^2SI^2}{2ml^3}t^2\right)}.
\]
此式适用于 $0\le t\le t_c$，其中
$t_c=\sqrt{2ml^3/(\mu_0N^2SI^2)}$。磁能减少恰等于两质点动能；$x\downarrow0$ 时理想模型的长螺线管近似失效，不能继续外推。

图~\ref{advA4:quality:A4-magnetostatics:a3} 中，长度变化改变自感 $L(x)$，理想超导闭回路保持磁链而不是电流；两端以相反速度运动，长度坐标的等效质量须由两质点动能相加得到。
\end{advsolution}

\subsection{磁介质与磁标势}

% CompletenessEM:a4-magnetic-medium
\noindent\textbf{从磁偶极矩密度到束缚电流。}
磁化强度 $\vect M$ 是单位体积的磁偶极矩。令
$R=|\vect r-\vect r'|$，体元的偶极矢势叠加为
\[
 \vect A_M(\vect r)=\frac{\mu_0}{4\pi}\int_V
 \vect M(\vect r')\times\nabla'\frac1R\dd V'.
\]
由 $\nabla'\times(\vect M/R)
=\nabla'(1/R)\times\vect M+(\nabla'\times\vect M)/R$，
并把旋度体积分转为边界积分，得到
\[
 \vect A_M=\frac{\mu_0}{4\pi}\left[
 \int_V\frac{\nabla'\times\vect M}{R}\dd V'
 +\oint_{\partial V}\frac{\vect M\times\vect n}{R}\dd S'\right].
\]
与电流的矢势比较，便识别出
\[
 \boxed{\vect j_M=\nabla\times\vect M,\qquad
 \vect K_M=\vect M\times\vect n}.
\]
若界面两侧都有磁化，$\vect n$ 从介质1指向介质2，则净面流为
$(\vect M_1-\vect M_2)\times\vect n$。
均匀磁化体内部 $\vect j_M=0$，并不意味着没有磁场，边界面流仍然存在。

\begin{qualityfigure}
\begin{tikzpicture}[>=Latex,font=\small,scale=.9]
 \fill[PaleBlue](-1.6,-1.1)rectangle(1.3,1.35);
 \draw[thick](1.3,-1.1)--(1.3,1.35);
 \foreach\x in{-.8,.2}{\draw[->,DeepBlue,thick](\x,-.65)--(\x,.7);}
 \node[left]at(-.8,.1){$\vect M$};
 \draw[->](1.3,.9)--(2.45,.9)node[right]{$\vect n=\vect e_x$};
 \node[draw,circle,inner sep=1pt]at(1.3,-.35){$\times$};
 \node[right]at(1.7,-.35){$\vect K_M=M\vect e_y$（入纸面）};
 \draw[->](-1.8,-1.65)--(-.65,-1.65)node[right]{$x$};
 \draw[->](-1.8,-1.65)--(-1.8,-.6)node[above]{$z$};
\end{tikzpicture}
\figcap{右手坐标的 $xz$ 截面：$+y$ 指向纸内，因此 $\vect e_z\times\vect e_x=\vect e_y$ 的面流画作叉号。\label{complete:em:bound-current}}
\end{qualityfigure}

图~\ref{complete:em:bound-current} 把 $\vect M$、外法线和面流三个方向同时固定，
可用于检查下一例磁化球的面流符号。定义
$\vect H=\vect B/\mu_0-\vect M$ 后，静磁 Maxwell 方程成为
\[
 \nabla\times\vect H=\vect j_f,\qquad
 \nabla\cdot\vect H=-\nabla\cdot\vect M.
\]
在无自由电流的单连通区域才能全局写 $\vect H=-\nabla\psi_M$；于是
\[
 \rho_M=-\nabla\cdot\vect M,\qquad
 \sigma_M=(\vect M_1-\vect M_2)\cdot\vect n,
\]
\[
 \nabla^2\psi_M=-\rho_M,\qquad
 \psi_M(\vect r)=\frac1{4\pi}\left[
 \int\frac{\rho_M}{R}\dd V'+\oint\frac{\sigma_M}{R}\dd S'\right].
\]
这些磁荷只是计算 $\vect H$ 的等效源；实际 $\vect B=\mu_0(\vect H+\vect M)$
仍满足 $\nabla\cdot\vect B=0$，不能把等效磁荷理解为真实磁单极。
若有自由电流，可先求其真空形式的场
\[
 \vect H_f=\frac1{4\pi}\int
 \frac{\vect j_f(\vect r')\times(\vect r-\vect r')}{R^3}\dd V',
 \qquad \vect H=\vect H_f-\nabla\psi_M,
\]
再联立磁化关系与边界条件确定 $\vect M$。

对准静态、无磁滞且可逆的介质，固定材料参数时增量能量为
$\dd w=\vect H\cdot\dd\vect B$。
若本构关系是通过原点的线性各向同性关系
$\vect M=\chi_m\vect H$、$\vect B=\mu\vect H$，从零场积分得到
\[
 w=\int_0^{\vect B}\vect H(\vect B')\cdot\dd\vect B'
 =\frac12\vect B\cdot\vect H.
\]
有永久磁化、磁滞或色散时，不能不加条件地套用这一简单能量式。



\begin{advexamplebox}{均匀磁化球的内外磁场}\label{advA4:ex:uniform-magnetized-sphere}
半径 $R$ 的球体具有均匀磁化强度 $\vect M=M\vect e_z$。求内外 $\vect H,\vect B$。

\begin{center}\begin{minipage}{.96\linewidth}\centering
\begin{tikzpicture}[>=Latex,font=\small,scale=.85,every node/.style={fill=none}]

\draw[thick] (0,0) circle (1.2);\foreach \x in {-.5,0,.5}{\draw[DeepBlue,->] (\x,-.6)--(\x,.65);}\node[above] at (0,1.5){$\boldsymbol M=M\boldsymbol e_z$};\draw[->] (0,0)--(1.2,0) node[midway,below]{$R$};\node[right] at (1.7,.4){球外无自由电流};

\end{tikzpicture}
\figcap{题图：均匀磁化球、磁化轴与半径，内外场均待求。}\label{advA4:quality:A4-magnetostatics:q4}
\end{minipage}\end{center}
\end{advexamplebox}
\global\AdvHideHeadingtrue
\begin{advsolution}
\begin{center}\begin{minipage}{.96\linewidth}\centering
% EMAdvFigureHeading
\raggedright{\bfseries\color{DeepBlue}解答}\par\smallskip\centering

\begin{tikzpicture}[>=Latex,font=\small,scale=.85,every node/.style={fill=none}]

\draw[thick] (0,0) circle (1.2);\foreach \x in {-.5,0,.5}{\draw[DeepBlue,->] (\x,-.6)--(\x,.65);}\node[above] at (0,1.5){$\boldsymbol M=M\boldsymbol e_z$};\draw[->] (0,0)--(1.2,0) node[midway,below]{$R$};\node[right] at (1.7,.4){球外无自由电流};

\end{tikzpicture}
\figcap{解图：均匀磁化球、磁化轴与半径，内外场均待求。}\label{advA4:quality:A4-magnetostatics:a4}
\end{minipage}\end{center}

体磁荷 $-\nabla\cdot\vect M=0$，表面磁荷
$\sigma_M=M\cos\theta$。轴对称标势取
\[
 \psi_{\rm in}=Ar\cos\theta,\qquad
 \psi_{\rm out}=\frac{B\cos\theta}{r^2}.
\]
切向 $H$ 连续，法向满足
$H_{n,\rm out}-H_{n,\rm in}=\sigma_M$。代入得到
\[
 \boxed{\vect H_{\rm in}=-\frac13\vect M},
 \qquad
 \boxed{\vect B_{\rm in}=\frac{2\mu_0}{3}\vect M}.
\]
球外等效于位于球心的磁偶极子
\[
 \vect m=\frac{4\pi R^3}{3}\vect M,
\]
所以
\[
 \boxed{\vect B_{\rm out}=\frac{\mu_0}{4\pi r^3}
 [3(\vect m\cdot\vect e_r)\vect e_r-\vect m]}.
\]
球外是真空，另有 $\boldsymbol H_{\rm out}=\boldsymbol B_{\rm out}/\mu_0$。在 $r=R$ 检查 $B_n$ 连续、$H_t$ 连续，边界条件满足。

图~\ref{advA4:quality:A4-magnetostatics:a4} 中，均匀磁化使体束缚电流为零、面束缚电流为 $\boldsymbol M\times\boldsymbol n$；内外标势的连续和法向跃变决定退磁场，不能把 $H$ 与 $B/\mu_0$ 在球内等同。
\end{advsolution}

\begin{advexamplebox}{自由电流圈的磁偶极层与磁标势}\label{advA4:ex:current-loop-scalar-potential}
证明载流 $I$ 的闭合线圈在无电流的单连通子区域内可由磁标势描述，并等效于跨任意张面 $S$ 的磁偶极层；说明跨张面的支路跳跃。

\begin{center}\begin{minipage}{.96\linewidth}\centering
\begin{tikzpicture}[>=Latex,font=\small,scale=.85,every node/.style={fill=none}]

\fill[DeepBlue!8] (0,0) ellipse (1.4 and .65);\draw[thick] (0,0) ellipse (1.4 and .65);\draw[DeepBlue,->] (1.4,0) arc[start angle=0,end angle=65,x radius=1.4,y radius=.65];\draw[->] (0,0)--(0,1.5) node[above]{$\boldsymbol n$};\node at (.3,-.2){张面 $S$};\node[right] at (1.7,0){边界电流 $I$};

\end{tikzpicture}
\figcap{题图：闭合电流圈与任取的跨面，张面法向和边界绕行方向按右手规则关联。}\label{advA4:quality:A4-magnetostatics:q5}
\end{minipage}\end{center}
\end{advexamplebox}
\global\AdvHideHeadingtrue
\begin{advsolution}
\begin{center}\begin{minipage}{.96\linewidth}\centering
% EMAdvFigureHeading
\raggedright{\bfseries\color{DeepBlue}解答}\par\smallskip\centering

\begin{tikzpicture}[>=Latex,font=\small,scale=.85,every node/.style={fill=none}]

\fill[DeepBlue!8] (0,0) ellipse (1.4 and .65);\draw[thick] (0,0) ellipse (1.4 and .65);\draw[DeepBlue,->] (1.4,0) arc[start angle=0,end angle=65,x radius=1.4,y radius=.65];\draw[->] (0,0)--(0,1.5) node[above]{$\boldsymbol n$};\node at (.3,-.2){张面 $S$};\node[right] at (1.7,0){边界电流 $I$};

\end{tikzpicture}
\figcap{解图：闭合电流圈与任取的跨面，张面法向和边界绕行方向按右手规则关联。}\label{advA4:quality:A4-magnetostatics:a5}
\end{minipage}\end{center}

张面法线与正电流方向满足右手关系，其有向立体角定义为
\[
 \Omega(\vect r)=\int_S
 \frac{(\vect r-\vect r')\cdot\dd\vect S'}{|\vect r-\vect r'|^3}.
\]
% CompletenessEM:a4-loop-stokes
下面把这个梯度具体化。记 $f=1/|\vect r-\vect r'|$，则
$\nabla'f=(\vect r-\vect r')/R^3$，且对核有 $\partial_i f=-\partial'_i f$。
在场点不在张面上时，$\nabla'^2f=0$。用分量形式的 Stokes 公式，
\[
 \begin{aligned}
 \left[\oint_{\partial S}\dd\vect l'\times\nabla'f\right]_i
 &=\int_S\left(n_j\partial'_i\partial'_jf
       -n_i\nabla'^2f\right)\dd S'\\
 &=\int_S n_j\partial'_i\partial'_jf\dd S'
 =-\partial_i\int_S\vect n\cdot\nabla'f\dd S'
 =-\partial_i\Omega.
 \end{aligned}
\]
第一行可由 $\varepsilon_{ijk}\varepsilon_{j\ell m}
=\delta_{im}\delta_{k\ell}-\delta_{i\ell}\delta_{km}$ 展开验证，
这也明确了源点导数与场点导数之间的负号。

\begin{qualityfigure}
\begin{tikzpicture}[>=Latex,font=\small,scale=.9]
 \fill[PaleBlue](0,0)ellipse(1.65 and .6);
 \draw[thick](0,0)ellipse(1.65 and .6);
 \draw[DeepBlue,->,thick](1.65,0)arc[start angle=0,end angle=65,x radius=1.65,y radius=.6];
 \node[right]at(1.65,.25){$I,\ \partial S$};
 \draw[->](0,-.95)--(0,1.6)node[above]{$\vect n$};
 \fill(0,.85)circle(1.7pt)node[right]{$h>0$};
 \fill(0,-.75)circle(1.7pt)node[right]{$h<0$};
 \draw[dashed](-.65,0)ellipse(.65 and .23);
 \draw[<->](-.65,-1.25)--(0,-1.25)node[midway,below]{$a_0$};
\end{tikzpicture}
\figcap{张面局部用半径 $a_0$ 的平面片核对支路跳跃；正法向与边界电流按右手规则关联。\label{complete:em:loop-cut}}
\end{qualityfigure}

图~\ref{complete:em:loop-cut} 中，令 $h$ 为场点相对局部平面片的有符号高度，
其立体角贡献为
\[
 \Omega_{\rm loc}(h)=2\pi\int_0^{a_0}
 \frac{hs\dd s}{(s^2+h^2)^{3/2}}
 =2\pi\left(\operatorname{sgn}h-\frac{h}{\sqrt{h^2+a_0^2}}\right).
\]
$h\to0^\pm$ 时分别为 $\pm2\pi$；远离这个小面片的部分连续，
所以跨割面的跳跃恰为 $4\pi$。这里不是说磁场在任意张面上产生真实面流奇点，
而是对两侧支路分别求梯度。综上得到


\[
 -\nabla\left(\frac I{4\pi}\Omega\right)
 =\frac I{4\pi}\oint_{\partial S}
 \frac{\dd\vect l'\times(\vect r-\vect r')}{R^3}
 =\vect H.
\]
因此
\[
 \boxed{\psi_H=\frac I{4\pi}\Omega},
 \qquad \vect B=-\mu_0\nabla\psi_H.
\]
按正法线方向穿过割面，$\Omega$ 从 $-2\pi$ 跳到 $+2\pi$，故选定支上的 $\psi_H$ 跳变 $+I$；沿正向闭合回路跟踪连续支，势累计下降 $I$，与环路积分
$\oint\vect H\cdot\dd\vect l=I$。面上单位面积磁偶极矩取
$I\vect n$ 就产生同一势，所以线圈等效于磁偶极层。标势是多值的，只有在割去张面后的单连通区域内才能单值使用。

图~\ref{advA4:quality:A4-magnetostatics:a5} 中，改变张面只改变标势所用的支路面，真实磁场由边界电流决定；跨面势函数可有跳跃，而绕过电流圈的环路不允许全局单值标势。
\end{advsolution}

\subsection{静磁边值问题}

无自由电流区域内 $\vect H=-\nabla\psi$ 且 $\nabla^2\psi=0$。理想超导体完全抗磁，内部 $\vect B=0$，表面给出
$\partial_n\psi=0$ 的 Neumann 条件。

\begin{advexamplebox}{反向平行电流、二维磁偶极与超导圆柱}\label{advA4:ex:parallel-currents-cylinder}
两条无限长 $z$ 向导线分别位于 $x=+d/2$ 和 $x=-d/2$，相应载流为 $+I$ 和 $-I$。求磁标势；取 $p=Id$、$d\to0$ 的二维偶极极限。再把该偶极置于半径 $R$ 的无限长超导圆柱空腔轴线上，求腔内平衡场。

\begin{center}\begin{minipage}{.96\linewidth}\centering
\begin{tikzpicture}[>=Latex,font=\small,scale=.85,every node/.style={fill=none}]

\draw[thick] (0,0) circle (1.4);\node at (-.4,0){$\otimes$};\node at (.4,0){$\odot$};\node[above] at (-.4,.3){$-I$};\node[above] at (.4,.3){$+I$};\draw[<->] (-.4,-.8)--(.4,-.8) node[midway,below]{$d$};\node[above] at (0,1.7){超导圆柱空腔横截面 $R$};

\end{tikzpicture}
\figcap{题图：两反向长电流及其二维偶极极限，柱面边界在横截面内取圆。}\label{advA4:quality:A4-magnetostatics:q6}
\end{minipage}\end{center}
\end{advexamplebox}
\global\AdvHideHeadingtrue
\begin{advsolution}
\begin{center}\begin{minipage}{.96\linewidth}\centering
% EMAdvFigureHeading
\raggedright{\bfseries\color{DeepBlue}解答}\par\smallskip\centering

\begin{tikzpicture}[>=Latex,font=\small,scale=.85,every node/.style={fill=none}]

\draw[thick] (0,0) circle (1.4);\node at (-.4,0){$\otimes$};\node at (.4,0){$\odot$};\node[above] at (-.4,.3){$-I$};\node[above] at (.4,.3){$+I$};\draw[<->] (-.4,-.8)--(.4,-.8) node[midway,below]{$d$};\node[above] at (0,1.7){超导圆柱空腔横截面 $R$};

\end{tikzpicture}
\figcap{解图：两反向长电流及其二维偶极极限，柱面边界在横截面内取圆。}\label{advA4:quality:A4-magnetostatics:a6}
\end{minipage}\end{center}

在不跨过割线的区域，单根 $+I$ 导线的标势是
$-I\arg(z-z_0)/(2\pi)$。因此在共同选定的割线之外，成对电流的标势为
\[
 \boxed{\psi=-\frac I{2\pi}
 [\arg(x-d/2+\ii y)-\arg(x+d/2+\ii y)]}.
\]
小 $d$ 展开得到
\[
 \boxed{\psi_0=-\frac p{2\pi}\frac y{x^2+y^2}
 =-\frac p{2\pi r}\sin\phi}.
\]
超导边界要求 $\partial_r\psi|_{R}=0$。加入在圆柱内正则的 $l=1$ 调和项，
\[
 \boxed{\psi=-\frac p{2\pi}\left(\frac{\sin\phi}{r}
 +\frac{r\sin\phi}{R^2}\right)}.
\]
于是 $\vect H=-\nabla\psi$，直角坐标分量为
\[
 H_x=-\frac{pxy}{\pi r^4},
 \qquad
 H_y=\frac p{2\pi}\left(\frac{x^2-y^2}{r^4}+\frac1{R^2}\right),
 \qquad \vect B=\mu_0\vect H.
\]
在 $r=R$ 直接计算 $H_r=0$；例如自由空间中 $x=0,y\ne0$ 处有 $H_y=-p/(2\pi y^2)$，也与右侧 $+I$、左侧 $-I$ 的 Biot--Savart 场相符。$R\to\infty$ 时附加均匀场消失，退回自由空间偶极。

图~\ref{advA4:quality:A4-magnetostatics:a6} 中，两导线的标势必须使用相同支路约定，取 $d\to0$ 时保持 $Id$ 固定；超导边界要求法向磁场为零，不能只规定标势本身为零。
\end{advsolution}

\begin{advexamplebox}{圆柱面电流与三圆柱相交体}\label{advA4:ex:steinmetz-currents}
无限薄圆柱面半径为 $R$，分别载
$\vect K=i_0\vect e_\phi$ 与
$\vect K=i_0\cos\phi\,\vect e_z$，求内外磁场。

三个互相垂直、半径均为 $R$ 的圆柱公共区域组成三圆柱相交体。其三类边界面上分别给定
\[
 \vect K_1=-i_0\cos\phi_y\vect e_x,\quad
 \vect K_2=i_0\cos\phi_x\vect e_y,\quad
 \vect K_3=i_0\vect e_{\phi_z}.
\]
证明电流连续并求远场。

\begin{center}
\begin{tikzpicture}[scale=.9]
 \draw[axis](-2,-1)--(2.4,1.2)node[right]{$x$};
 \draw[axis](2,-1)--(-2.4,1.2)node[left]{$y$};
 \draw[axis](0,-1.5)--(0,2.1)node[above]{$z$};
 \draw[DeepBlue] (0,0)ellipse(1.8 and .9);
 \draw[vector,Gold!80!black](.7,.8)--(-.3,.3);\node[anchor=west]at(.8,1.55){$\vect K_1\parallel-\vect e_x$};\draw[guide](.8,1.4)--(.2,.55);
 \draw[vector,DeepBlue](-.7,.8)--(-1.7,1.3)node[above]{$\vect K_2\parallel\vect e_y$};
 \draw[vector,TextGray](1.8,0)--(1.8,.65)node[right]{$\vect K_3$ 切向};
 \node at(0,-2){局部方向示意，非三圆柱相交体的投影边界};
\end{tikzpicture}
\figcap{题图：原圆柱与三圆柱相交边界的切向电流几何保留。}\label{advA4:quality:A4-magnetostatics:q7}

\end{center}
\end{advexamplebox}
\global\AdvHideHeadingtrue
\begin{advsolution}
\begin{center}\begin{minipage}{.96\linewidth}\centering
% EMAdvFigureHeading
\raggedright{\bfseries\color{DeepBlue}解答}\par\smallskip\centering

\begin{tikzpicture}[scale=.9]
 \draw[axis](-2,-1)--(2.4,1.2)node[right]{$x$};
 \draw[axis](2,-1)--(-2.4,1.2)node[left]{$y$};
 \draw[axis](0,-1.5)--(0,2.1)node[above]{$z$};
 \draw[DeepBlue] (0,0)ellipse(1.8 and .9);
 \draw[vector,Gold!80!black](.7,.8)--(-.3,.3);\node[anchor=west]at(.8,1.55){$\vect K_1\parallel-\vect e_x$};\draw[guide](.8,1.4)--(.2,.55);
 \draw[vector,DeepBlue](-.7,.8)--(-1.7,1.3)node[above]{$\vect K_2\parallel\vect e_y$};
 \draw[vector,TextGray](1.8,0)--(1.8,.65)node[right]{$\vect K_3$ 切向};
 \node at(0,-2){局部方向示意，非三圆柱相交体的投影边界};
\end{tikzpicture}
\figcap{解图：原圆柱与三圆柱相交边界的切向电流几何保留。}\label{advA4:quality:A4-magnetostatics:a7}
\end{minipage}\end{center}

环形面电流正是均匀磁化 $\vect M=i_0\vect e_z$ 的束缚面电流
$\vect K=\vect M\times\vect n$，所以
\[
 \boxed{\vect B_{\rm in}=\mu_0i_0\vect e_z,\qquad
 \vect B_{\rm out}=0}.
\]
对 $K_z=i_0\cos\phi$，取 $A_z$：
\[
 A_{z,\rm in}=ar\cos\phi,\qquad
 A_{z,\rm out}=\frac b r\cos\phi.
\]
$B_r$ 连续和
$B_{\phi,\rm out}-B_{\phi,\rm in}=\mu_0K_z$ 给出
$a=\mu_0i_0/2$、$b=aR^2$。于是
\[
 \boxed{\vect B_{\rm in}=-\frac{\mu_0i_0}{2}\vect e_y},
\]
\[
 \boxed{B_{r,\rm out}=-\frac{\mu_0i_0R^2}{2r^2}\sin\phi,\quad
 B_{\phi,\rm out}=\frac{\mu_0i_0R^2}{2r^2}\cos\phi}.
\]

三圆柱相交体的关键是三组电流并非彼此拼凑。取整个体积均匀磁化
$\vect M=i_0\vect e_z$，则任意边界面上
$\vect K=\vect M\times\vect n$。在沿 $x$ 的圆柱面，
$\vect n=\cos\phi_y\vect e_y+\sin\phi_y\vect e_z$，得到
$\vect K_1=-i_0\cos\phi_y\vect e_x$；沿 $y$、$z$ 的圆柱面同理恰得 $\vect K_2,\vect K_3$。因此接缝两侧是同一个连续面流的限制，电荷不在接缝堆积，连续性自动满足。

远场只保留总磁偶极矩
\[
 \vect m=\int\vect M\dd V=i_0V_S\vect e_z.
\]
% CompletenessEM:a4-triple-volume
为计算 $V_S$，把三圆柱公共区域写成
\[
 x^2+y^2\leq R^2,\quad x^2+z^2\leq R^2,\quad y^2+z^2\leq R^2.
\]
先限于第一卦限，再取其中的 $x\geq y$ 部分；其余由三坐标变号和 $x,y$ 互换生成，
总体积是这一部分的 $16$ 倍。对固定的 $x,y$，后两个不等式给出
$0\leq z\leq\sqrt{R^2-x^2}$；投影区域如图~\ref{complete:em:triple-projection}，
$y$ 的上界为 $\min(x,\sqrt{R^2-x^2})$，在 $x=R/\sqrt2$ 处分段。

\begin{qualityfigure}
\begin{tikzpicture}[>=Latex,font=\small,scale=1.05]
 \fill[PaleBlue](0,0)--(1.6971,1.6971)
 arc[start angle=45,end angle=0,radius=2.4]--cycle;
 \draw[->](0,0)--(2.85,0)node[right]{$x$};
 \draw[->](0,0)--(0,2.8)node[above]{$y$};
 \draw[DeepBlue,thick](0,0)--(1.6971,1.6971);
 \draw[DeepBlue,thick](1.6971,1.6971)arc[start angle=45,end angle=0,radius=2.4];
 \draw[TextGray](0,2.4)arc[start angle=90,end angle=45,radius=2.4];
 \draw[dashed](1.6971,0)--(1.6971,1.6971);
 \node[below left]at(1.6971,-.05){$R/\sqrt2$};\node[below right]at(2.4,-.05){$R$};
 \node[above left]at(1,1){$y=x$};
 \node[above right]at(2.05,1.35){$x^2+y^2=R^2$};
 \node at(1.25,.5){$D$};
\end{tikzpicture}
\figcap{积分域 $D$ 是第一卦限且 $x\geq y$ 的 $xy$ 投影；每点竖向高度为 $\sqrt{R^2-x^2}$。\label{complete:em:triple-projection}}
\end{qualityfigure}

因此不用预先记忆相交体积，也可直接积分：
\[
 \begin{aligned}
 V_S&=16\left[\int_0^{R/\sqrt2}x\sqrt{R^2-x^2}\dd x
 +\int_{R/\sqrt2}^{R}(R^2-x^2)\dd x\right]\\
 &=16R^3\left[\frac13-\frac1{6\sqrt2}
 +\frac23-\frac5{6\sqrt2}\right]
 =8(2-\sqrt2)R^3.
 \end{aligned}
\]
第一段来自 $0\leq y\leq x$，第二段来自 $0\leq y\leq\sqrt{R^2-x^2}$；
这就是数值系数中出现 $\sqrt2$ 的几何来源。


三个正交圆柱公共区域的体积为
$V_S=8(2-\sqrt2)R^3$，故
\[
 \boxed{\vect m=8(2-\sqrt2)i_0R^3\vect e_z},
\]
\[
 \boxed{\vect B(\vect r)\sim\frac{\mu_0}{4\pi r^3}
 [3(\vect m\cdot\vect e_r)\vect e_r-\vect m]}
 \qquad(r\gg R).
\]
图中只示三类面流的局部方向，不把椭圆当作三圆柱相交体的边界。各面电流均由同一公式
$\vect M\times\vect n$ 决定，所以方向关系不是装饰性示意，而直接参与连续性证明。

图~\ref{advA4:quality:A4-magnetostatics:a7} 中，面电流只能沿实际边界切向流动；接缝处各面流入、流出量须连续，远区偶极矩要对三类面电流共同积分，不能任意把某一完整圆柱场截断。
\end{advsolution}
